\chapter{1958年江苏卷}

\section{解答题}

\begin{problem}
\begin{enumerate}
\item
解方程 \(\lg(x+7)-\lg(x+1)=\lg(x+3)\).

\item
解方程 \(\sin 10x=\frac{1}{2}\).

\item
已知 \(\tan\frac{\theta}{2}=t\)，求证
\(\tan\theta=\frac{2t}{1-t^{2}}\)，
\(\sin\theta=\frac{2t}{1+t^{2}}\)，
\(\cos\theta=\frac{1-t^{2}}{1+t^{2}}\).
\end{enumerate}
\end{problem}

\begin{solution}
\begin{enumerate}
\item
由对数的定义，必须有 \(x+7>0\)、\(x+1>0\)、\(x+3>0\)，所以 \(x>-1\).

在此条件下，原方程等价于 \(\lg\frac{x+7}{x+1}=\lg(x+3)\). 由于对数函数在正数范围内是单调函数，所以 \(x+7=(x+1)(x+3)\)，即 \(x^{2}+3x-4=(x-1)(x+4)=0\). 因此 \(x=1\) 或 \(x=-4\). 其中 \(x=-4\) 不满足 \(x>-1\)，舍去；代入可知 \(x=1\) 满足原方程，所以原方程的解为 \(x=1\).

\item
由 \(\sin u=\frac{1}{2}\) 的通解，令 \(u=10x\)，得 \(10x=\frac{\pi}{6}+2k\pi\) 或 \(10x=\frac{5\pi}{6}+2k\pi\)，其中 \(k\in\Z\). 分别除以 \(10\)，得到 \(x=\frac{\pi}{60}+\frac{k\pi}{5}\) 或 \(x=\frac{\pi}{12}+\frac{k\pi}{5}\)，其中 \(k\in\Z\). 将这两族数值分别代入可知都满足 \(\sin 10x=\frac12\)，故它们构成原方程的全部解.

\item
因为 \(\tan\frac{\theta}{2}=t\)，所以 \(\cos\frac{\theta}{2}\ne0\). 由正切的定义和平方关系，\(1=\sin^{2}\frac{\theta}{2}+\cos^{2}\frac{\theta}{2}=\left(1+t^{2}\right)\cos^{2}\frac{\theta}{2}\)，从而 \(\cos^{2}\frac{\theta}{2}=\frac{1}{1+t^{2}}\). 利用二倍角公式，\(\sin\theta=2\sin\frac{\theta}{2}\cos\frac{\theta}{2}=2\tan\frac{\theta}{2}\cos^{2}\frac{\theta}{2}=\frac{2t}{1+t^{2}}\)，以及 \(\cos\theta=\cos^{2}\frac{\theta}{2}-\sin^{2}\frac{\theta}{2}=\left(1-t^{2}\right)\cos^{2}\frac{\theta}{2}=\frac{1-t^{2}}{1+t^{2}}\).
当 \(t\ne\pm1\) 时，\(\cos\theta\ne0\)，因此 \(\tan\theta=\frac{\sin\theta}{\cos\theta}=\frac{\dfrac{2t}{1+t^{2}}}{\dfrac{1-t^{2}}{1+t^{2}}}=\frac{2t}{1-t^{2}}\).
当 \(t=\pm1\) 时，\(\cos\theta=0\)，此时 \(\tan\theta\) 无定义；所以正切公式应理解为在 \(\tan\theta\) 有定义的条件下成立.
\end{enumerate}
\end{solution}

\begin{problem}
解方程 \(\sin x+\sqrt{3}\cos x=2\).
\end{problem}

\begin{solution}
利用和角公式，\(\sin x+\sqrt{3}\cos x=2\left(\sin x\cos\frac{\pi}{3}+\cos x\sin\frac{\pi}{3}\right)=2\sin\left(x+\frac{\pi}{3}\right)\). 因此原方程等价于 \(\sin\left(x+\frac{\pi}{3}\right)=1\). 正弦值等于 \(1\) 当且仅当 \(x+\frac{\pi}{3}=\frac{\pi}{2}+2k\pi\)，其中 \(k\in\Z\)，所以 \(x=\frac{\pi}{6}+2k\pi\)，其中 \(k\in\Z\).

\end{solution}

\begin{problem}
已知正六棱锥底边长为 \(a\)，高为 \(h\)，求侧面积.
\end{problem}

\begin{solution}
\solutionfigure{\bitmapfigure[max width=0.33\linewidth]{img/1958/jiangsu/q3_solution_fig1.png}}

设棱锥顶点为 \(A\)，正六边形底面的中心为 \(O\)，取一条底边 \(BC\)，并设 \(H\) 为 \(BC\) 的中点. 正六边形的外接圆半径等于边长，所以 \(OB=a\)；又因中心到边的垂线平分该边，故 \(BH=\frac{a}{2}\)，且 \(OH\perp BC\). 在直角三角形 \(OBH\) 中，\(OH=\sqrt{OB^{2}-BH^{2}}=\sqrt{a^{2}-\frac{a^{2}}{4}}=\frac{\sqrt{3}}{2}a\). 正棱锥的高垂直于底面，故 \(AO=h\)，从而 \(AO\perp BC\)，且 \(\triangle AOH\) 为直角三角形. 由于 \(BC\) 垂直于相交直线 \(AO\)、\(OH\)，所以 \(BC\perp\) 平面 \(AOH\)，特别地 \(AH\perp BC\)，即 \(AH\) 是侧面三角形 \(ABC\) 的高. 因此 \(AH=\sqrt{AO^{2}+OH^{2}}=\sqrt{h^{2}+\frac{3}{4}a^{2}}\). 每一个侧面三角形的面积为 \(S_{\triangle ABC}=\frac{1}{2}a\cdot AH=\frac{a}{2}\sqrt{h^{2}+\frac{3}{4}a^{2}}\). 正六棱锥有六个全等的侧面，因此侧面积为 \(S=6S_{\triangle ABC}=3a\sqrt{h^{2}+\frac{3}{4}a^{2}}\).

\end{solution}

\begin{problem}
已知长方体三条棱长之和为 \(a+b+c=13\)，对角线为 \(\sqrt{65}\)，求全面积.
\end{problem}

\begin{solution}
长方体的体对角线满足 \(a^{2}+b^{2}+c^{2}=\left(\sqrt{65}\right)^{2}=65\). 长方体的全面积为 \(S=2\left(ab+bc+ca\right)\). 又因为 \(\left(a+b+c\right)^{2}=a^{2}+b^{2}+c^{2}+2\left(ab+bc+ca\right)\)，所以 \(S=\left(a+b+c\right)^{2}-\left(a^{2}+b^{2}+c^{2}\right)=13^{2}-65=104\).
故长方体的全面积为 \(104\).

\end{solution}

\begin{problem}
已知 \(AB\) 为直径，\(DC\) 为弦，四边形 \(ABCD\) 内接于 \(\odot O\)，并记 \(BC=a\)，\(CD=b\)，\(DA=c\). 求证：

\begin{enumerate}
\item \(AC\cdot DB=AB\cdot DC+AD\cdot CB\).

\item \(AB\) 为方程 \(x^{3}-\left(a^{2}+b^{2}+c^{2}\right)x-2abc=0\) 的根.
\end{enumerate}
\end{problem}

\begin{solution}
\solutionfigure{\bitmapfigure[max width=0.34\linewidth]{img/1958/jiangsu/q5_solution_fig1.png}}

\begin{enumerate}
\item
在边 \(AC\) 上取点 \(E\)，使射线 \(DE\) 位于 \(\angle CDA\) 内，且满足 \(\angle ADE=\angle BDC\). 因为四边形 \(ABCD\) 内接于同一个圆，且 \(E\) 在 \(AC\) 上，所以 \(\angle DAE=\angle DAC=\angle DBC\). 结合 \(\angle ADE=\angle BDC\)，可得 \(\triangle ADE\sim\triangle BDC\)，于是 \(\frac{AD}{BD}=\frac{AE}{BC}\)，即 \(AD\cdot CB=AE\cdot BD\).

另一方面，\(\angle DCE=\angle DCA=\angle DBA\). 又因为射线 \(DB\) 位于 \(\angle CDA\) 内部，所以 \(\angle CDE=\angle CDA-\angle ADE=\left(\angle CDB+\angle BDA\right)-\angle BDC=\angle BDA\). 因此 \(\triangle CDE\sim\triangle BDA\)，从而 \(\frac{CD}{BD}=\frac{CE}{BA}\)，即 \(AB\cdot DC=CE\cdot BD\). 将上面两式相加，并利用 \(AE+CE=AC\)，得 \(AB\cdot DC+AD\cdot CB=\left(CE+AE\right)BD=AC\cdot DB\). 这就证明了 \(AC\cdot DB=AB\cdot DC+AD\cdot CB\).

\item
因为 \(AB\) 是圆的直径，所以由圆周角定理，\(\angle ACB=\frac{\pi}{2}\)，且 \(\angle ADB=\frac{\pi}{2}\). 在直角三角形 \(ABC\) 和 \(ADB\) 中，分别有 \(AC^{2}=AB^{2}-BC^{2}=AB^{2}-a^{2}\) 以及 \(DB^{2}=AB^{2}-AD^{2}=AB^{2}-c^{2}\). 由第一问，且 \(CD=b\)，可得 \(AC\cdot DB=AB\cdot b+ac\).
令 \(X=AB\)，由于长度为正数，\(X>0\)，且 \(AC\)，\(DB>0\)，于是
\[
\sqrt{X^{2}-a^{2}}\sqrt{X^{2}-c^{2}}=bX+ac
\]
两边平方，得到
\[
X^{4}-\left(a^{2}+c^{2}\right)X^{2}+a^{2}c^{2}
=b^{2}X^{2}+2abcX+a^{2}c^{2}
\]
移项并约去相同项，得
\[
X^{4}-\left(a^{2}+b^{2}+c^{2}\right)X^{2}-2abcX=0
\]
由于 \(X>0\)，可除以 \(X\)，所以
\[
X^{3}-\left(a^{2}+b^{2}+c^{2}\right)X-2abc=0
\]
将 \(X=AB\) 代回，即知 \(AB\) 是方程
\[
x^{3}-\left(a^{2}+b^{2}+c^{2}\right)x-2abc=0
\]
的根.
\end{enumerate}
\end{solution}
